\section{Prediction against an endogenous label}\label{app:pred}

This appendix holds the machinery behind Section~\ref{sec:scoring}: the
general result, the binary and continuous constructions, the calibration
table, the fixed-detection grid, the finite-sample distributions and the
robustness checks. In audit terms, every result below says the same thing:
a score that steers where the audit looks can validate well against the
findings the audit then records while knowing nothing about where the
conditions are.

\subsection{Setup}

Let $S$ be a risk score, $C \in \{0,1\}$ the underlying condition and
$F \in \{0,1\}$ the recorded finding, with no false positives. Let $P(C=1) = p$
and let detection depend on the score, $d(S) = P(F=1 \mid C=1, S)$. Assume
throughout that $S \perp C$: the score carries \emph{no} information about
where conditions are. Then $P(F=1 \mid S) = p\,d(S)$, and any apparent
performance of $S$ against $F$ is produced entirely by the dependence of
detection on the score. The assumption is deliberately extreme; it isolates
the mechanism. A real score with some genuine information inherits the same
effect on top of whatever genuine performance it has, and the two are
indistinguishable from the validation output alone.

\subsection{A binary score}

Let $S \sim \mathrm{Bernoulli}(\pi)$ independent of $C$, with detection $d_1$
when $S=1$ and $d_0$ when $S=0$. Write $\bar d = \pi d_1 + (1-\pi) d_0$ and
$\mu = P(F=1) = p\,\bar d$.

\begin{proposition}\label{prop:aucbin}
$\displaystyle \mathrm{AUC}(S, F) = \tfrac12 +
\frac{\pi\,p\,(1-\pi)\,(d_1 - d_0)}{2\,\mu\,(1-\mu)}$, while
$\mathrm{AUC}(S, C) = \tfrac12$ exactly.
\end{proposition}

The second statement is immediate from $S \perp C$. At the prespecified values
$p = \nBinP{}$, $\pi = \nBinPi{}$, $d_0 = \nBinDzero{}$ and $d_1 = \nBinDone{}$,
the closed form gives $\mathrm{AUC}(S,F) = \nBinAucExact{}$; a Monte Carlo run
at $N = 10^6$ returns \nBinAucMC{}.

\paragraph{A labeling trap worth naming.} The score in
Proposition~\ref{prop:aucbin} is a generic binary risk indicator drawn
independently, not the lagged finding itself. If the score \emph{is} the
previous period's recorded finding, $S = F_{t-1}$, in the stationary regime of
Proposition~\ref{prop:markov}, then
%
\[
  \mathrm{AUC}(F_{t-1}, F_t) \;=\; \tfrac12 + \tfrac12\,p\,(d_1 - d_0)
  \;\leq\; \frac{1+p}{2},
\]
%
which at $p = \nBinP{}$ is bounded by \nBinStationaryBound{}, far below the
\nBinAucExact{} available to a generic score. Within the observation-only
process the lagged finding is a weak predictor of the record, because it is
constrained by the same prevalence that limits the record. That is a statement
about this model and not about lagged findings in general: where the
underlying condition itself persists, the lagged finding carries information
about that condition too, and the bound does not apply.

\subsection{A continuous score}

Let $S$ be continuous with a logistic detection curve $d(S)$ of steepness
$\alpha$, again independent of $C$. As $\alpha$ grows, detection becomes a
sharper function of the score and $\mathrm{AUC}(S,F)$ can be driven toward one
while $\mathrm{AUC}(S,C)$ remains exactly one-half. At $\alpha = 40$ the exact
value by quadrature is $\mathrm{AUC}(S,F) = \nContAucExact{}$, with a Monte
Carlo value of \nContAucMC{}, against a Monte Carlo $\mathrm{AUC}(S,C) =
\nContAucC{}$.

\subsection{Calibration}

\begin{proposition}\label{prop:calib}
Let $m(S) = p\,d(S)$. Then $E[F \mid m] = m$: the model is perfectly calibrated
for the recorded finding. But $E[C \mid m] = p$ for every value of $m$: it
carries no information about the underlying condition at any point of its
range.
\end{proposition}

\begin{proof}
$E[F \mid S] = P(F=1 \mid S) = p\,d(S) = m(S)$, and $m$ is a function of $S$, so
$E[F \mid m] = m$. Since $S \perp C$, $E[C \mid S] = p$, hence $E[C \mid m] = p$.
\end{proof}

Table~\ref{tab:calib} shows this numerically. Across deciles of the score, mean
predictions run from \nCalibLo{} to \nCalibHi{}, a span of roughly
\nCalibSpan{}-fold, and the recorded finding rate tracks them closely, while
the underlying problem rate sits between \nCalibLatentLo{} and
\nCalibLatentHi{} in every decile. A calibration plot against an endogenously
observed label is a statement about that label.

\subsection{Fixed average detection}

The constructions above can be dismissed as saying only that more attention
produces more findings. This one cannot. The quantity held fixed is $\bar d$,
the average probability of detecting an existing condition across the
portfolio; it is not a measure of audit hours, staffing or cost. Fix $\bar d$
and set
%
\[
  d_0 = \bar d - \pi k, \qquad d_1 = \bar d + (1-\pi)\,k .
\]
%
Then average detection $\pi d_1 + (1-\pi) d_0 = \bar d$ is invariant in $k$, and
so is the recorded prevalence $\mu = p\,\bar d$. The parameter $k$ moves
detection from unflagged cases to flagged ones without changing the total.

At the prespecified values $p = \nBudgetP{}$, $\pi = \nBudgetPi{}$ and
$\bar d = \nBudgetDetection{}$, evaluated at \nBudgetPoints{} points on
$k \in [0, 0.8]$, six quantities behave as the construction requires and were
checked at every point. Average detection stays at \nBudgetDetection{} and
recorded prevalence at \nBudgetPrevalence{}. The AUC against the recorded
finding rises monotonically from \nBudgetAucFstart{} to \nBudgetAucFend{}, with
the closed form and the direct computation agreeing throughout. The AUC
against the underlying condition is exactly \nBudgetAucC{} at every point. The
Brier score against the recorded finding falls monotonically from
\nBudgetBrierFstart{} to \nBudgetBrierFend{}, and the Brier score against the
underlying condition rises monotonically from \nBudgetBrierCstart{} to
\nBudgetBrierCend{}. Table~\ref{tab:budget} gives the grid. The configuration
was fixed in advance and no parameter search was performed.

\subsection{Robustness of the predictive results}

\paragraph{Sample size.} Figure~\ref{fig:finite} gives the sampling
distributions at portfolio sizes an audit organization actually has. At
\nFiniteNsmall{} programs the middle 95 percent of $\mathrm{AUC}(S,F)$ runs
from \nFiniteAucFlo{} to \nFiniteAucFhi{}, a range a reviewer would read as
clear evidence of a working model, and $\mathrm{AUC}(S,C)$ brackets one-half
at every sample size tested, with \nFiniteReps{} replications per
configuration.

\paragraph{A fitted model rather than an oracle.} The scores above are oracles:
the exact functional form is known. Fitting a conventional flexible learner
by cross-validation instead recovers the oracle closely and inherits its
property: correlation with the oracle between \nFitCorrLo{} and \nFitCorrHi{},
$\mathrm{AUC}(S,F) = \nFitAucF{}$ against an oracle value of
\nFitAucFOracle{}, and $\mathrm{AUC}(S,C) = \nFitAucC{}$. Nothing about the
phenomenon requires the analyst to know the true detection function; an
ordinary model discovers it from the data.

\paragraph{False positives.} Relaxing the no-false-positive assumption, so that
a finding can be recorded where no condition exists at rates $a_0$ and $a_1$,
attenuates the recorded risk ratio but does not overturn the qualitative
conclusion: with $a_0 = a_1 = \nFPrate{}$, $d_0 = \nFPdzero{}$ and
$d_1 = \nFPdone{}$, \nFPshare{}~percent of recorded findings are false positives,
$\mathrm{AUC}(S,F) = \nFPaucF{}$ and $\mathrm{AUC}(S,C) = \nFPaucC{}$.

\begin{table}[htbp]
\centering
\caption{A fixed detection-capacity budget, every fourth grid point of 21.}
\label{tab:budget}
\small
\begin{tabular}{ccccccccc}
\toprule
$k$ & $d_0$ & $d_1$ & \shortstack{average\\detection} & \shortstack{recorded\\rate}
& AUC vs $F$ & AUC vs $C$ & Brier vs $F$ & Brier vs $C$ \\
\midrule
0.00 & 0.400 & 0.400 & 0.400 & 0.080 & 0.5000 & 0.5000 & 0.0736 & 0.1744 \\
0.16 & 0.360 & 0.520 & 0.400 & 0.080 & 0.5408 & 0.5000 & 0.0734 & 0.1746 \\
0.32 & 0.320 & 0.640 & 0.400 & 0.080 & 0.5815 & 0.5000 & 0.0728 & 0.1752 \\
0.48 & 0.280 & 0.760 & 0.400 & 0.080 & 0.6223 & 0.5000 & 0.0719 & 0.1761 \\
0.64 & 0.240 & 0.880 & 0.400 & 0.080 & 0.6630 & 0.5000 & 0.0705 & 0.1775 \\
0.80 & 0.200 & 1.000 & 0.400 & 0.080 & 0.7038 & 0.5000 & 0.0688 & 0.1792 \\
\bottomrule
\end{tabular}
\tablenotes{What is held fixed is $\bar d$, the average probability of detecting an
existing deficiency; it is not a measure of audit hours or cost.
Detection is set by $d_0 = \bar d - \pi k$ and $d_1 = \bar d + (1-\pi)k$, with
$p = 0.20$, $\pi = 0.25$ and
$\bar d = 0.40$, all prespecified. Average detection and
recorded prevalence are invariant in $k$ by construction; the table shows they are.
Exact calculation.}
\end{table}

\begin{table}[htbp]
\centering
\caption{Calibrated for the record, uninformative about the problem. Deciles of
the oracle score $m(S) = p\,d(S)$ at $\alpha = 20$. Predicted values track the
recorded finding rate across a span of roughly 94-fold,
while the underlying problem rate is flat at $p = 0.20$ in every
decile. Simulation, $N = 10^6$.}
\label{tab:calib}
\small
\begin{tabular}{cccc}
\toprule
Decile of the score & Mean prediction & Recorded finding rate & Underlying problem rate \\
\midrule
1 & 0.0020 & 0.0019 & 0.2004 \\
2 & 0.0020 & 0.0022 & 0.1996 \\
3 & 0.0020 & 0.0020 & 0.2006 \\
4 & 0.0020 & 0.0022 & 0.2009 \\
5 & 0.0022 & 0.0022 & 0.1983 \\
6 & 0.0035 & 0.0036 & 0.2011 \\
7 & 0.0127 & 0.0132 & 0.2024 \\
8 & 0.0573 & 0.0566 & 0.1997 \\
9 & 0.1417 & 0.1423 & 0.1998 \\
10 & 0.1872 & 0.1852 & 0.1983 \\
\bottomrule
\end{tabular}
\end{table}


\begin{figure}[htbp]
\centering
\includegraphics[width=0.6\textwidth]{exB3_finite.pdf}
\caption{\capFinite}
\label{fig:finite}
\end{figure}
